\paperhead
{A Fuzzy Logic Controller for an Assistive-Care Flat: Design, Genetic Tuning\\
of Its Membership Functions, and an Optimiser Comparison on CEC'2005}
{%
We design a fuzzy logic controller that manages heating, cooling and
lighting in one room of an assistive-care flat for a resident with limited
mobility. Its 54 rules are written in plain linguistic terms so that a carer
can check them, and the controller is implemented in both Python and MATLAB,
which produce the same outputs. A genetic algorithm then tunes the controller,
reducing its normalised test error by 32.6\% and clearly outperforming random search given
the same effort. The tuning came at a cost, however: three of five runs
scrambled the order of the fuzzy sets, so some rules no longer meant what they
said, and four of five left parts of the input range with no output at all; no
run avoided both problems. Finally, a genetic algorithm, particle swarm
optimisation and simulated annealing are compared on two standard benchmark
functions. None wins everywhere: particle swarm is best in two dimensions, the
genetic algorithm is the most reliable in ten, while the ranking of
simulated annealing depends on the statistic used.%
}
{fuzzy logic control, Mamdani inference, membership function tuning, genetic
algorithms, particle swarm optimisation, simulated annealing, CEC'2005,
assistive technology}

